\documentclass[12pt,a4paper]{article}

% --- PACKAGES ---
\usepackage[utf8]{inputenc}
\usepackage[french]{babel}
\usepackage[T1]{fontenc}
\usepackage{geometry}
\geometry{top=2.5cm, bottom=2.5cm, left=2cm, right=2cm}
\usepackage{graphicx}
\usepackage{enumitem}
\usepackage{xcolor}
\usepackage{titlesec}
\usepackage{hyperref}
\usepackage{fancyhdr}
\usepackage{lastpage}
\usepackage{tcolorbox}
\usepackage{amssymb}
\usepackage{lmodern}
\usepackage{tikz} % Pour la déco graphique
\usetikzlibrary{shapes,arrows,positioning}

% --- PALETTE DE COULEURS BÉNIN (PROFESSIONNEL) ---
\definecolor{beningreen}{RGB}{0, 135, 81}
\definecolor{beninyellow}{RGB}{252, 209, 22}
\definecolor{beninred}{RGB}{232, 17, 45}
\definecolor{darkblue}{RGB}{20, 40, 80}
\definecolor{softgrey}{RGB}{240, 240, 240}

% --- CONFIGURATION DES TITRES ---
\titleformat{\section}{\Large\bfseries\color{beningreen}}{\thesection}{1em}{}[\titlerule]
\titleformat{\subsection}{\large\bfseries\color{darkblue}}{\thesubsection}{1em}{}
\titleformat{\subsubsection}{\normalsize\bfseries\color{beninred!80!black}}{\thesubsubsection}{1em}{}

% --- CONFIGURATION DES EN-TÊTES ---
\pagestyle{fancy}
\fancyhf{}
\setlength{\headheight}{18pt}
\lhead{\small \textit{Bénin Geo-Watch : Intelligence Artificielle \& Géopolitique}}
\rhead{\small \textbf{Équipe G10}}
\cfoot{\small Page \thepage\ sur \pageref{LastPage}}

\hypersetup{colorlinks=true, linkcolor=darkblue, urlcolor=beningreen}

\begin{document}

% ============================================================
% PAGE DE GARDE GRAPHIQUE
% ============================================================
\begin{titlepage}
    \begin{tikzpicture}[remember picture, overlay]
        \fill[softgrey] (current page.south west) rectangle (current page.north east);
        \fill[beningreen] (current page.north west) rectangle ([xshift=1cm]current page.south west);
        \node[anchor=north east, xshift=-1cm, yshift=-1cm] at (current page.north east) {\Large \textbf{HACKATHON 2026}};
    \end{tikzpicture}

    \centering
    \vspace*{2cm}
    
    {\Huge \bfseries \color{darkblue} BÉNIN GEO-WATCH \par}
    \vspace{0.5cm}
    {\Large \color{beningreen!70!black} Système Intégré de Surveillance et de Prédiction Géopolitique \par}
    
    \vspace{1.5cm}
    \begin{tikzpicture}
        \draw[beningreen, line width=2pt] (0,0) -- (10,0);
        \draw[beninyellow, line width=2pt] (0,-0.2) -- (10,-0.2);
        \draw[beninred, line width=2pt] (0,-0.4) -- (10,-0.4);
    \end{tikzpicture}
    
    \vspace{1.5cm}
    \begin{tcolorbox}[colback=white, colframe=darkblue, arc=5mm, boxrule=1pt, width=0.8\textwidth]
        \centering
        \large \textbf{RAPPORT FINAL DÉTAILLÉ (A-Z)} \\
        \small \textit{Analyse, Modélisation, RAG et Déploiement}
    \end{tcolorbox}
    
    \vspace{2cm}
    \begin{minipage}{0.5\textwidth}
        \begin{flushleft} \large
            \emph{Membres de l'Équipe G10 :}\\
            \rule{5cm}{0.5pt} \\
            \textbf{Ahamadah Marie Cynthia} \\ \small Data Analyst \\
            \vspace{0.2cm}
            \textbf{OLOUKPONA Bedel} \\ \small Data Science \\
            \vspace{0.2cm}
            \textbf{ADISSO Nelxie} \\ \small ML Engineer \\
            \vspace{0.2cm}
            \textbf{Nolhan-Nathan Luzolo ADOHO} \\ \small ML Engineer
        \end{flushleft}
    \end{minipage}
    \hfill
    \begin{minipage}{0.45\textwidth}
        \begin{flushright} \large
            \textbf{Livrables Clés} \\
            \small
            Dashboard Streamlit App \\
            Pipeline RAG Mistral \\
            Forecasting J+7 \\
            Clustering GDELT
        \end{flushright}
    \end{minipage}
    
    \vfill
    {\large \today}
\end{titlepage}

\newpage
\tableofcontents
\newpage

% ============================================================
\section{Vision Stratégique et Objectifs}
% ============================================================
Le projet \textbf{Bénin Geo-Watch} répond à l'impératif de souveraineté numérique et informationnelle. Dans un monde saturé de données, la capacité à extraire des \textbf{signaux faibles} est un avantage stratégique.

\subsection{Les trois piliers du projet}
\begin{enumerate}
    \item \textbf{Observation :} Cartographie départementale des tensions (12 départements couverts).
    \item \textbf{Anticipation :} Prédiction de la stabilité via Machine Learning.
    \item \textbf{Interaction :} Dialogue avec la connaissance via un assistant RAG.
\end{enumerate}

% ============================================================
\section{Data Engineering : Du Big Data au Smart Data}
% ============================================================
Le socle du projet repose sur la base \textbf{GDELT (Global Database of Events, Language and Tone)}.

\subsection{Ingestion et Nettoyage (Notebooks)}
Nous avons traité un dataset de plusieurs milliers d'événements. Le nettoyage a inclus :
\begin{itemize}
    \item \textbf{Filtrage CAMEO :} Identification des acteurs (Bénin, Nigeria, France, etc.).
    \item \textbf{Géocodage :} Extraction des coordonnées pour les 12 départements (Alibori, Atacora, Borgou, etc.).
    \item \textbf{Calcul d'Indices :} Création d'un \textit{Indice de Stabilité} composite normalisé sur une échelle de 0 à 100.
\end{itemize}

\begin{tcolorbox}[colback=beningreen!5, colframe=beningreen, title=Focus : Normalisation]
    L'échelle de Goldstein (-10 à +10) a été transformée en un score intuitif où \textbf{50 représente la neutralité}. Ce travail a été crucial pour rendre le dashboard compréhensible par un non-expert.
\end{tcolorbox}

% ============================================================
\section{Modélisation Prédictive et Clustering}
% ============================================================

\subsection{Clustering KMeans (Non-supervisé)}
Nous avons identifié 4 clusters géopolitiques distincts basés sur le comportement médiatique :
\begin{itemize}
    \item \textbf{Cluster 0 (Crise) :} Fort volume, Tone très négatif, Goldstein bas.
    \item \textbf{Cluster 2 (Coopération) :} Événements diplomatiques à impact positif.
\end{itemize}

\subsection{Forecasting RandomForest}
Pour la prédiction temporelle, nous utilisons un \textit{Random Forest Regressor}.
\subsubsection{Performance du modèle}
Le modèle atteint des métriques de précision satisfaisantes (\textbf{MAE faible}) grâce à l'utilisation de \textbf{Lag Features} (décalages temporels J-1 à J-7) et de moyennes mobiles, capturant ainsi la saisonnalité des tensions.

% ============================================================
\section{Architecture RAG : L'Intelligence Augmentée}
% ============================================================
L'Équipe G10 a déployé un pipeline de \textit{Retrieval-Augmented Generation} ultra-performant.

\subsection{Stack Technique RAG}
\begin{itemize}
    \item \textbf{Embeddings :} Modèle \texttt{mistral-embed} (Vecteurs de haute dimension).
    \item \textbf{Vector Store :} \textbf{FAISS} (\textit{Facebook AI Similarity Search}) pour une recherche instantanée parmi 6 353 articles.
    \item \textbf{Inférence :} \texttt{mistral-small-latest} pour la synthèse textuelle.
\end{itemize}

\begin{tcolorbox}[colback=beninred!5, colframe=beninred, title=Défis du RAG]
    Le principal défi était la \textbf{quantité absurde de données}. La solution a été la création d'un index vectoriel persisté, évitant de re-calculer les embeddings à chaque requête, réduisant la latence de 90\%.
\end{tcolorbox}

% ============================================================
\section{Développement du Dashboard (UI/UX)}
% ============================================================
Le dashboard a été conçu pour être un outil de "Command \& Control".

\subsection{Design Système}
\begin{itemize}
    \item \textbf{Thème Sombre Premium :} Réduction de la fatigue visuelle et aspect "Centre d'Opérations".
    \item \textbf{Visualisations Plotly :} Graphiques interactifs (Area charts, Mapbox) pour une exploration granulaire.
    \item \textbf{Simulateur What-If :} Permet à l'utilisateur de simuler un scénario (ex: chute brutale du Goldstein) et de voir comment le modèle ML le classifierait en temps réel.
\end{itemize}

% ============================================================
\section{Difficultés et Gestion de Projet}
% ============================================================
\begin{enumerate}
    \item \textbf{Partage PowerBI :} L'impossibilité de partager facilement l'analyse nous a poussé à migrer vers Streamlit. Cette contrainte s'est transformée en opportunité (intégration native de l'IA).
    \item \textbf{Travail en Équipe :} La coordination entre le nettoyage des données et le développement de l'interface a été facilitée par l'utilisation de \texttt{joblib} pour la sérialisation des modèles.
\end{enumerate}

% ============================================================
\section{Livrables et Accès}
% ============================================================
\begin{itemize}
    \item \textbf{Dashboard :} \href{https://isheeroo.streamlit.app/}{isheeroo.streamlit.app}
    \item \textbf{Code Source :} Dépôt GitHub avec structure \texttt{uv}.
    \item \textbf{Rapport :} Ce présent document LaTeX.
\end{itemize}

% ============================================================
\section{Conclusion}
% ============================================================
Bénin Geo-Watch est la preuve que la Data Science peut éclairer les zones d'ombre de la géopolitique. Notre équipe livre aujourd'hui un outil robuste, prédictif et prêt pour une utilisation stratégique.

\vfill
\begin{center}
    \textit{Équipe G10 -- Hackathon iSHEERO 2026} \\
    \textbf{Fait avec rigueur et passion pour le futur du Bénin.}
\end{center}

\end{document}
